\scanpage{60}
\begin{center}4 - 34 -\hfill n\textsuperscript{o} 392\end{center}

entier sur $A$, il n'y a pas d'autre idéal que $0$ au-dessus de l'idéal $0$ de $A$, (chap. V, § 2, n\textsuperscript{o} 1, cor. 1 de la prop. 1), donc il revient au même de dire que $\underline{a}\ne0$ ou que $\underline{b}\ne0$.

(i) En vertu de ce qui précède, on peut se borner au cas où $R$ est un $B$-module de torsion. Si $\underline{b}$ est contenu dans un idéal premier $\underline{P}\in P(B)$, $\underline{a}$ est contenu dans $\underline{P}\cap A=\underline{p}$, qui est de hauteur $1$. Inversement, si $\underline{a}$ est contenu dans un idéal premier $\underline{p}\in P(A)$, il existe un idéal premier $\underline{P}$ de $B$ qui contient $\underline{b}$ et est au-dessus de $\underline{p}$, (chap. V, § 2, n\textsuperscript{o} 1, cor. 2 du th. 1). L'assertion (i) résulte de ces remarques et du n\textsuperscript{o} 4, prop. 9.

(ii) Pour tout $B$-module de torsion de type fini $R$, posons $\varphi(R)=\unsure{\chi|\gamma}_A(R_{[A]})$ ; il est clair que, (pour les $B$-modules de torsion de type fini), $\varphi$ vérifie les conditions 1) et 2) de la prop. 11 du n\textsuperscript{o} 5, (compte tenu de (i)). Il existe par suite un homomorphisme $\unsure{\vartheta|\theta}:D(B)\longrightarrow D(A)$ tel que $\varphi(R)=\unsure{\vartheta|\theta}(\unsure{\chi|\gamma}_B(R))$ pour tout $B$-module de torsion $R$ de type fini. L'homomorphisme $\unsure{\vartheta|\theta}$ est déterminé par sa valeur pour tout $B$-module de la forme $B/\underline{P}$ où $\underline{P}\in P(B)$, puisque $\unsure{\chi|\gamma}_B(B/\underline{P})=\underline{P}$. Or, pour tout idéal premier $\underline{q}\ne\underline{p}=\underline{P}\cap A$ dans $P(A)$, on a $\underline{q}\not\subset\underline{P}$, donc $(B/\underline{P})_{\underline{q}}=0$. D'autre part, si on pose $S=A-\underline{p}$, $\underline{P}.S^{-1}B$ est un idéal maximal de $S^{-1}B$ et $(B/\underline{P})_{\underline{p}}=S^{-1}B/\underline{P}.S^{-1}B$ est isomorphe au corps des fractions de $B/\underline{P}$, c'est-à-dire au corps résiduel de $v_{\underline{P}}$ ; sa longueur en tant que $A_{\underline{p}}$-module est donc $f_{\underline{P}/\underline{p}}$ ; ce qui prouve que $\unsure{\vartheta|\theta}=N$, (n\textsuperscript{o} 5, déf. 4).

(iii) Si $T$ est le sous-module de torsion de $R$, $T_{[A]}$ est le sous-module de torsion de $R_{[A]}$, et $(R/T)_{[A]}=R_{[A]}/T_{[A]}$ ; pour prouver (21), on peut donc se borner au cas où $R$ est sans torsion. Alors $R$ est identifié à un sous-$B$-module de $R_{(L)}$, et contient une base $(e_i)_{1\le i\le m}$ de $R_{(L)}$ sur $L$. Si $(b_j)_{1\le j\le n}$ est une base de $L$ sur $K$ formée d'éléments de $B$, les $b_je_i$ constituent une base de $R_{(L)}$ sur $K$, formée d'éléments de $R$, d'où (21). Soit d'autre part $M$ un sous-$B$-module

\margin{59}
